\chapter*{Abstract}
\addcontentsline{toc}{chapter}{Abstract}

Modern cloud-native applications deployed on Kubernetes face an evolving DDoS threat
landscape that existing perimeter-focused defenses are ill-equipped to address. Current
solutions --- whether signature-based filtering, statistical anomaly detection, or
machine learning classifiers operating on individual network flows --- share a fundamental
limitation: they treat network traffic as a collection of independent data points, failing
to model the relational and temporal structure of microservice communication.

This project presents \sys{}, an intelligent, Kubernetes-native DDoS defense framework
that addresses this gap through three integrated innovations. First, we leverage
eBPF-based network observability via Cilium and Hubble to collect real-time,
fine-grained flow telemetry from within the Kubernetes data plane, capturing traffic
between pods and services with negligible overhead. Second, we construct dynamic
communication graphs from this telemetry and apply graph-based anomaly detection using
Graph Autoencoders and GraphSAGE to identify attack patterns in the relational structure
of microservice interactions. Third, we translate graph anomaly scores into adaptive,
graduated rate-limiting responses enforced via XDP/TC eBPF programs and dynamic
CiliumNetworkPolicy updates, preserving legitimate traffic while suppressing attack flows.

Experimental evaluation on a Kubernetes testbed running the Online Boutique benchmark
application demonstrates that \sys{} effectively detects volumetric, application-layer,
and low-and-slow DDoS attacks with lower false positive rates than conventional
flow-based classifiers, while adaptive mitigation reduces attack impact by [X]\% compared
with static rate limiting. This work contributes a reproducible, open-source prototype
and a labelled dataset of Kubernetes DDoS scenarios, providing a foundation for future
research in cloud-native network security.

\bigskip
\noindent \textbf{Keywords:} eBPF, Kubernetes, DDoS detection, Graph Neural Networks,
Graph Autoencoder, GraphSAGE, microservices, network security, anomaly detection,
adaptive rate limiting.
